\section{Optimasi Parameter pada Sistem Klasik}
Dalam pemodelan sistem makroskopik berskala besar, pencarian titik ekuilibrium pada algoritma kuantum merupakan hal yang sangat rentang terjadi fenomena \textit{noise} atau derau, sehingga memerlukan metode optimasi parameter yang stabil untuk menavigasi kompleksitas kuantum secara lebih kokoh. \textit{Noise} tersebut merupakan bagian dari realitas komputasi pada lingkungan fisik yang mendistorsi sinyal informasi secara signifikan sehingga menyulitkan proses kalibrasi model. Metode evaluasi gradien klasik yang sering digunakan pada skala industri sering kali memicu ketidakstabilan numerik karena operasi tersebut justru dapat meningkatkan derau pengukuran \cite{spall_implementation_1998}.

\subsection{Metode Pergeseran Parameter sebagai Estimator Ekspektasi}
Dalam setiap proses optimasi berparameter pada model komputasi kuantum, perhitungan nilai gradien pada dasarnya memang sangat diperlukan untuk mengarahkan pembaruan parameter menuju konfigurasi fungsi minimum. Metode konvensional seperti (\textit{finite difference}) sangat bergantung pada limit pergeseran infinitesimal yang sering digunakan pada sistem klasik nyatanya sangat sulit untuk diimplementasikan secara langsung pada arsitektur sistem fisik kuantum. Dikarenakan sistem komputasi fisik riil yang memiliki varians pengukuran persisten, operasi pembagian dengan nilai pergeseran yang sangat kecil akan meningkatkan derau statistik secara signifikan sehingga menyebabkan ketidakstabilan perhitungan energi \cite{wierichs_general_2022}. Di sisi lain, Metode Pergeseran Parameter (\textit{Parameter Shift Rule}) berhasil melenyapkan hambatan fundamental ini dengan mengeksploitasi struktur polinomial trigonometri pada fungsi objektif sistem sirkuit kuantum berparameter \cite{mitarai_quantum_2018}.

Secara matematis, fungsi objektif $E(\theta)$ pada sirkuit kuantum berparameter diperlakukan sebagai deret Fourier terbatas sehingga turunannya dapat dinyatakan secara eksak dengan mengevaluasi kombinasi linier fungsi tersebut \cite{wierichs_general_2022}. Untuk sistem komputasi dengan spektrum frekuensi tunggal, turunan analitik terhadap parameter sudut $\theta$ dapat diekspresikan melalui pergeseran fase simetris sebesar $\pm \pi/2$ sehingga dapat memunculkan sinyal gradien tanpa terdistorsi oleh derau pengukuran \cite{mitarai_quantum_2018}. Nilai yang diekstraksi melalui metode pergeseran ini merupakan representasi turunan analitis murni orde pertama. Selanjutnya, Evaluasi diskrit yang sudah diperoleh dapat memberikan estimator gradien tanpa bias (\textit{unbiased estimator}) sehingga pembaruan parameter untuk energi yang akurat dapat diformulasikan melalui perumusan analitik
\begin{equation}
\label{eq:parameter_shift_rule}
\frac{\partial E}{\partial \theta} = \frac{1}{2} \left[ E\left(\theta + \frac{\pi}{2}\right) - E\left(\theta - \frac{\pi}{2}\right) \right].
\end{equation}

\subsection{Optimasi \textit{Gradient Descent}}
Algoritma penurunan gradien atau yang biasa disebut (\textit{Gradient Descent}) merupakan salah satu metode optimasi komputasional yang paling fundamental dalam menavigasi ruang parameter suatu fungsi objektif. Proses pembaruan parameter pada setiap langkah iterasi ke-$t$ diatur sedemikian rupa dengan menggunakan formulasi matematis yang melibatkan vektor gradien $\nabla L(\theta_t)$ dan laju pembelajaran $\eta$ dengan bentuk
\begin{equation}
\label{eq:gradient_descent}
\theta_{t+1} = \theta_t - \eta \nabla L(\theta_t).
\end{equation}
Parameter skalar $\eta$ atau \textit{learning rate} berfungsi untuk mengontrol ukuran langkah, di mana nilai ekstrem dapat memicu divergensi numerik sementara nilai yang terlalu kecil akan memperlambat konvergensi sistem secara signifikan sehingga memakan daya komputasi yang berlebih. Dalam lanskap parameter berdimensi sangat tinggi, keberhasilan algoritma ini sangat terbebani oleh kompleksitas estimasi gradien yang akurat, sehingga sering kali rentan terjebak dalam jebakan \textit{local minima} yang merugikan \cite{mcclean_barren_2018}. Selain itu, algoritma ini mewajibkan eksekusi evaluasi minimal sebanyak dua kali pada setiap tahapan iterasi yang sedang berlangsung.

Salah satu keterbatasan dari optimasi \textit{Gradient Descent} pada arsitektur komputasi fisik adalah tingginya beban kompleksitas yang tumbuh secara linear terhadap dimensi vektor parameter $\theta$ \cite{spall_implementation_1998}. Kompleksitas komputasi metode \textit{Parameter Shift Rule} yang dibutuhkan untuk mengekstraksi satu vektor gradien analitik sangat bergantung terhadap jumlah parameter. Untuk mengekstraksi satu vektor gradien analitik menggunakan \textit{Parameter Shift Rule}, arsitektur sistem wajib melakukan eksekusi evaluasi dua kali pada setiap tahapan iterasi sehingga akan memakan daya komputasi dua kali lipat dari jumlah iterasi untuk evaluasi gradien \cite{wierichs_general_2022}. Ketika model optimasi tersebut diimplementasikan pada simulasi portofolio finansial berskala makro, jumlah variabel parameter $p$ akan mengalami kenaikan eksponensial seiring dengan peningkatan jumlah aset dan kedalaman sirkuit kuantum yang dioperasikan. Keterbatasan operasional tersebut memunculkan peralihan paradigma optimasi menuju metode stokastik yang mampu mereduksi hambatan komputasi tanpa mengorbankan kualitas maupun stabilitas konvergensi sistem, salah satunya adalah metode SPSA \cite{spall_implementation_1998}.

\subsection{Optimasi SPSA (\textit{Simultaneous Perturbation Stochastic Approximation})}
Algoritma \textit{Simultaneous Perturbation Stochastic Approximation} (SPSA) merupakan metode optimasi yang mereduksi kompleksitas evaluasi gradien analitik dari skala $\mathcal{O}(p)$ menjadi $\mathcal{O}(1)$ \cite{spall_implementation_1998}. Meskipun mereduksi kompleksitas komputasi, algoritma ini berhasil mencapai efisiensi asimtotik karena hanya membutuhkan dua kali evaluasi fungsi objektif untuk mengestimasi seluruh komponen vektor gradien berdimensi tinggi secara simultan. Estimasi pada algoritma ini beroperasi dengan menginjeksikan perturbasi stokastik secara paralel pada seluruh parameter sistem, tidak seperti optimasi metode \textit{finite difference} yang harus melalui pergeseran sekuensial secara individual. Reduksi kompleksitas komputasi ini merupakan hal yang sangat krusial bagi optimasi pada arsitektur sistem kuantum dikarenakan pada komputasi kuantum, setiap eksekusi fungsi tujuan mengonsumsi sumber daya komputasi yang masif. Oleh karena itu, implementasi SPSA menjadi salah satu metode optimasi yang menjanjikan bagi pencarian solusi optimal pada arsitektur dengan masalah berskala eksponensial yang sebelumnya tidak mungkin didekati menggunakan metode gradien konvensional \cite{wang_variational_2025}.

Pada setiap siklus iterasi ke-$k$, algoritma SPSA akan membangkitkan sebuah vektor perturbasi stokastik $\Delta_k$ di mana seluruh elemennya $(k=1, 2, 3, \dots)$ ditarik dari distribusi Rademacher yang memiliki nilai $\pm 1$. Komponen ke-$i$ dari estimasi vektor gradien $\hat{g}_{k}(\theta_k)$ diformulasikan sebagai selisih evaluasi objektif yang dinormalisasi terhadap besaran perturbasi pada komponen spesifik tersebut dengan bentuk
\begin{equation}
\label{eq:estimasi_gradien_spsa}
\hat{g}_{k,i}(\theta_k) = \frac{L(\theta_k + c_k \Delta_k) - L(\theta_k - c_k \Delta_k)}{2 c_k \Delta_{k,i}}
\end{equation}
dimana $L(\cdot)$ merupakan fungsi objektif yang akan diminimalkan. Agar lintasan konvergensi stokastik dapat tercapai secara stabil, arsitektur SPSA menggunakan mekanisme peluruhan asimtotik bagi laju pembelajaran ($a_k$) maupun besaran perturbasi ($c_k$) seiring berjalannya tahapan iterasi. Dinamika hiperparameter tersebut diformulasikan agar mematuhi syarat konvergensi Robbins-Monro standar melalui pengaturan polinomial yang dapat dinyatakan dalam bentuk persamaan
\begin{equation}
\label{eq:peluruhan_hiperparameter}
a_k = \frac{a}{(A + k)^\alpha}, \quad c_k = \frac{c}{k^\gamma}.
\end{equation}
Peluruhan fungsi $c_k$ berfungsi untuk mengecilkan galat pemotongan secara bertahap, sementara pelemahan $a_k$ bertujuan untuk meredam fluktuasi varians yang timbul akibat injeksi perturbasi simultan \cite{spall_implementation_1998}.
